% figures/w04-track-frame.svg — 두 오차가 회전변환 하나에서 나온다.
%
%   bash _tools/tikz2svg.sh figures/src/w04-track-frame.tex figures/w04-track-frame.svg
%
% 다리는 북에서 pi_p = 62 deg. 배는 다리 위 32 m 지점에서 우현으로 14 m.
% [rotate=28] 안에서 다리를 +x 로 두면 우현이 국소 -y 다 (90 - 62 = 28).
\documentclass[border=5pt]{standalone}
% gnc-style.tex — 이 볼트의 TikZ 그림이 공유하는 스타일.
%
% 저널 그림 규약을 스타일 하나로 굳혀 둔다. 그림마다 선 굵기나 화살촉을
% 다시 정하지 않는다 — 그것이 그림이 제각각으로 보이는 원인이다.
%
% 쓰는 법:  % gnc-style.tex — 이 볼트의 TikZ 그림이 공유하는 스타일.
%
% 저널 그림 규약을 스타일 하나로 굳혀 둔다. 그림마다 선 굵기나 화살촉을
% 다시 정하지 않는다 — 그것이 그림이 제각각으로 보이는 원인이다.
%
% 쓰는 법:  % gnc-style.tex — 이 볼트의 TikZ 그림이 공유하는 스타일.
%
% 저널 그림 규약을 스타일 하나로 굳혀 둔다. 그림마다 선 굵기나 화살촉을
% 다시 정하지 않는다 — 그것이 그림이 제각각으로 보이는 원인이다.
%
% 쓰는 법:  \input{gnc-style}  (figures/src/ 안에서)

\usepackage{tikz}
\usepackage{amsmath}
\usepackage{xcolor}
\usetikzlibrary{arrows.meta, positioning, calc, fit, backgrounds, decorations.pathmorphing}

% ---- 색 --------------------------------------------------------------------
% 색은 강조 하나에만 쓴다. 나머지는 검정.
\definecolor{ink}{HTML}{1A1A1A}
\definecolor{accent}{HTML}{7C3AED}
\definecolor{warn}{HTML}{B91C1C}
\definecolor{soft}{HTML}{666666}

% ---- 선과 화살촉 -----------------------------------------------------------
\tikzset{
  % 신호선. 화살촉은 작고 뾰족한 것 하나뿐이다.
  sig/.style   = {ink, line width=0.5pt, -{Stealth[length=4pt, width=3pt]}},
  wire/.style  = {ink, line width=0.5pt},
  % 강조 경로 — 파선으로 구분한다. 흑백 인쇄에서도 살아야 하기 때문이다.
  asig/.style  = {accent, line width=0.5pt, densely dashed,
                  -{Stealth[length=4pt, width=3pt]}},
  awire/.style = {accent, line width=0.5pt, densely dashed},
  % 블록. 흰 바탕, 직각, 검은 테두리.
  blk/.style   = {draw=ink, line width=0.55pt, fill=white,
                  minimum width=17mm, minimum height=8mm, inner sep=2pt, font=\small},
  % 합산점
  sum/.style   = {draw=ink, line width=0.55pt, fill=white, circle,
                  inner sep=0pt, minimum size=4.6mm, font=\scriptsize},
  asum/.style  = {draw=accent, line width=0.55pt, fill=white, circle,
                  inner sep=0pt, minimum size=4.6mm, font=\scriptsize},
  % 분기점
  dot/.style   = {ink, fill=ink, circle, inner sep=0pt, minimum size=1.5mm},
  % 신호 이름 — 선 위에 작게
  sname/.style = {font=\small, inner sep=1.5pt},
  % 그림 안의 짧은 설명. 그림 안의 글은 최소로 둔다
  note/.style  = {font=\scriptsize, soft, inner sep=1.5pt},
  anote/.style = {font=\scriptsize, accent, inner sep=1.5pt},
  % 축
  axis/.style  = {ink, line width=0.5pt},
  tick/.style  = {font=\scriptsize, soft},
}

% 캡션 한 줄 — 그림 아래 가는 선과 함께
\newcommand{\gnccaption}[2]{%
  \node[anchor=north west, text width=#1, align=left, font=\scriptsize, ink]
        at (current bounding box.south west) {\vspace{2pt}#2};}
  (figures/src/ 안에서)

\usepackage{tikz}
\usepackage{amsmath}
\usepackage{xcolor}
\usetikzlibrary{arrows.meta, positioning, calc, fit, backgrounds, decorations.pathmorphing}

% ---- 색 --------------------------------------------------------------------
% 색은 강조 하나에만 쓴다. 나머지는 검정.
\definecolor{ink}{HTML}{1A1A1A}
\definecolor{accent}{HTML}{7C3AED}
\definecolor{warn}{HTML}{B91C1C}
\definecolor{soft}{HTML}{666666}

% ---- 선과 화살촉 -----------------------------------------------------------
\tikzset{
  % 신호선. 화살촉은 작고 뾰족한 것 하나뿐이다.
  sig/.style   = {ink, line width=0.5pt, -{Stealth[length=4pt, width=3pt]}},
  wire/.style  = {ink, line width=0.5pt},
  % 강조 경로 — 파선으로 구분한다. 흑백 인쇄에서도 살아야 하기 때문이다.
  asig/.style  = {accent, line width=0.5pt, densely dashed,
                  -{Stealth[length=4pt, width=3pt]}},
  awire/.style = {accent, line width=0.5pt, densely dashed},
  % 블록. 흰 바탕, 직각, 검은 테두리.
  blk/.style   = {draw=ink, line width=0.55pt, fill=white,
                  minimum width=17mm, minimum height=8mm, inner sep=2pt, font=\small},
  % 합산점
  sum/.style   = {draw=ink, line width=0.55pt, fill=white, circle,
                  inner sep=0pt, minimum size=4.6mm, font=\scriptsize},
  asum/.style  = {draw=accent, line width=0.55pt, fill=white, circle,
                  inner sep=0pt, minimum size=4.6mm, font=\scriptsize},
  % 분기점
  dot/.style   = {ink, fill=ink, circle, inner sep=0pt, minimum size=1.5mm},
  % 신호 이름 — 선 위에 작게
  sname/.style = {font=\small, inner sep=1.5pt},
  % 그림 안의 짧은 설명. 그림 안의 글은 최소로 둔다
  note/.style  = {font=\scriptsize, soft, inner sep=1.5pt},
  anote/.style = {font=\scriptsize, accent, inner sep=1.5pt},
  % 축
  axis/.style  = {ink, line width=0.5pt},
  tick/.style  = {font=\scriptsize, soft},
}

% 캡션 한 줄 — 그림 아래 가는 선과 함께
\newcommand{\gnccaption}[2]{%
  \node[anchor=north west, text width=#1, align=left, font=\scriptsize, ink]
        at (current bounding box.south west) {\vspace{2pt}#2};}
  (figures/src/ 안에서)

\usepackage{tikz}
\usepackage{amsmath}
\usepackage{xcolor}
\usetikzlibrary{arrows.meta, positioning, calc, fit, backgrounds, decorations.pathmorphing}

% ---- 색 --------------------------------------------------------------------
% 색은 강조 하나에만 쓴다. 나머지는 검정.
\definecolor{ink}{HTML}{1A1A1A}
\definecolor{accent}{HTML}{7C3AED}
\definecolor{warn}{HTML}{B91C1C}
\definecolor{soft}{HTML}{666666}

% ---- 선과 화살촉 -----------------------------------------------------------
\tikzset{
  % 신호선. 화살촉은 작고 뾰족한 것 하나뿐이다.
  sig/.style   = {ink, line width=0.5pt, -{Stealth[length=4pt, width=3pt]}},
  wire/.style  = {ink, line width=0.5pt},
  % 강조 경로 — 파선으로 구분한다. 흑백 인쇄에서도 살아야 하기 때문이다.
  asig/.style  = {accent, line width=0.5pt, densely dashed,
                  -{Stealth[length=4pt, width=3pt]}},
  awire/.style = {accent, line width=0.5pt, densely dashed},
  % 블록. 흰 바탕, 직각, 검은 테두리.
  blk/.style   = {draw=ink, line width=0.55pt, fill=white,
                  minimum width=17mm, minimum height=8mm, inner sep=2pt, font=\small},
  % 합산점
  sum/.style   = {draw=ink, line width=0.55pt, fill=white, circle,
                  inner sep=0pt, minimum size=4.6mm, font=\scriptsize},
  asum/.style  = {draw=accent, line width=0.55pt, fill=white, circle,
                  inner sep=0pt, minimum size=4.6mm, font=\scriptsize},
  % 분기점
  dot/.style   = {ink, fill=ink, circle, inner sep=0pt, minimum size=1.5mm},
  % 신호 이름 — 선 위에 작게
  sname/.style = {font=\small, inner sep=1.5pt},
  % 그림 안의 짧은 설명. 그림 안의 글은 최소로 둔다
  note/.style  = {font=\scriptsize, soft, inner sep=1.5pt},
  anote/.style = {font=\scriptsize, accent, inner sep=1.5pt},
  % 축
  axis/.style  = {ink, line width=0.5pt},
  tick/.style  = {font=\scriptsize, soft},
}

% 캡션 한 줄 — 그림 아래 가는 선과 함께
\newcommand{\gnccaption}[2]{%
  \node[anchor=north west, text width=#1, align=left, font=\scriptsize, ink]
        at (current bounding box.south west) {\vspace{2pt}#2};}


\begin{document}
\begin{tikzpicture}[x=1mm, y=1mm]

\def\XE{51.2}      % x_e^{\,p} = 32 m at 1.6 mm/m
\def\YE{22.4}      % y_e^{\,p} = 14 m
\def\LG{104}       % 다리 65 m

% ---- NED --------------------------------------------------------------------
\begin{scope}[shift={(-4,54)}]
  \draw[axis, -{Stealth[length=4pt,width=3pt]}] (0,0) -- (0,10);
  \draw[axis, -{Stealth[length=4pt,width=3pt]}] (0,0) -- (10,0);
  \node[font=\small, anchor=east]  at (-0.8,9) {$x^n$};
  \node[font=\small, anchor=north] at (9,-0.8) {$y^n$};
  \node[note, anchor=north west] at (0,-3) {the frame that never moves};
\end{scope}

\begin{scope}[rotate=28]
  % 다리
  \draw[pathc, line width=0.9pt, -{Stealth[length=4.5pt,width=3.4pt]}] (0,0) -- (\LG,0);
  \filldraw[pathc, fill=white, line width=0.7pt] (0,0) circle (1.1);

  % 다리가 정하는 프레임 : t 는 다리를 따라, n 은 우현으로
  \draw[pathc, line width=0.6pt, densely dashed, -{Stealth[length=4pt,width=3pt]}]
        (0,0) -- (22,0);
  \draw[pathc, line width=0.6pt, densely dashed, -{Stealth[length=4pt,width=3pt]}]
        (0,0) -- (0,-22);
  \node[pathc, font=\small, anchor=south] at (17,1.2) {$\hat{\mathbf{t}}$};
  \node[pathc, font=\small, anchor=west]  at (1.5,-19) {$\hat{\mathbf{n}}$ (to starboard)};

  % 위치오차 벡터
  \draw[ink, line width=0.9pt, -{Stealth[length=4.5pt,width=3.4pt]}]
        (0,0) -- (\XE,-\YE);
  % 두 성분
  \draw[alng, line width=0.9pt, -{Stealth[length=4.5pt,width=3.4pt]}] (0,0) -- (\XE,0);
  \draw[errc, line width=0.9pt, -{Stealth[length=4.5pt,width=3.4pt]}]
        (\XE,0) -- (\XE,-\YE);
  \draw[errc, line width=0.4pt] (\XE,-2.4) -- ({\XE-2.4},-2.4) -- ({\XE-2.4},0);

  % 배 — 다리 방향으로 대충 놓는다 (이 그림의 주제가 아니다)
  \fill[ink] ({\XE+3.6},{-\YE+1.9}) -- ({\XE-3.4},{-\YE+2.6})
          -- ({\XE-1.4},{-\YE-3.4}) -- cycle;

  \node[pathc, font=\small, anchor=north east] at (-1.5,-1.5) {$\mathbf{p}_i^{\,n}$};
  \node[pathc, font=\small, anchor=north west] at ({\LG-16},-2) {$\mathbf{p}_{i+1}^{\,n}$};
  \node[alng, font=\small, anchor=south]  at ({0.5*\XE},1.5) {$x_e^{\,p}$};
  \node[errc, font=\small, anchor=west]   at ({\XE+2},{-0.55*\YE}) {$y_e^{\,p}$};
  \node[ink, font=\scriptsize, anchor=north east] at ({0.62*\XE},{-0.62*\YE-3.5})
       {the position error};
  \node[ink, font=\scriptsize, anchor=north] at ({\XE+2},{-\YE-4}) {the vessel};
\end{scope}

% pi_p
\draw[soft, line width=0.4pt, densely dashed] (0,0) -- (0,30);
\draw[soft, line width=0.5pt] (0,23) arc (90:28:23);
\node[soft, font=\small] at (7,26) {$\pi_p$};

% ===================== 오른쪽 : 대수 =======================================
\begin{scope}[shift={(122,0)}]
  \node[font=\small\bfseries, anchor=north west] at (0,54) {one rotation, two answers};

  \node[anchor=north west, align=left, text width=68mm, font=\scriptsize] at (0,48) {%
    The leg gives the frame its direction:};
  \node[anchor=north west, font=\small] at (0,41)
       {$\pi_p = \mathrm{atan2}(y_{i+1}^n-y_i^n,\ x_{i+1}^n-x_i^n)$};
  \node[anchor=north west, align=left, text width=68mm, font=\scriptsize] at (0,34) {%
    and the error, written in that frame, is};
  \node[anchor=north west, font=\small] at (0,27)
       {$\begin{bmatrix} x_e^{\,p} \\ y_e^{\,p} \end{bmatrix}
         = \mathbf{R}_p^{\,n}(\pi_p)^{\!\top}
           \begin{bmatrix} x^n-x_i^n \\ y^n-y_i^n \end{bmatrix}$};

  \node[anchor=north west, align=left, text width=68mm, font=\scriptsize] at (0,6) {%
    $x_e^{\,p}$ — how far \textbf{along} the leg the vessel has come.
    The switching test of §4-6 uses it.\\[3pt]
    $y_e^{\,p}$ — how far \textbf{to the side}. The guidance law drives this one to
    zero.\\[6pt]
    Both are projections of the same vector onto the same orthonormal pair, so
    $x_e^2 + y_e^2 = \lVert\mathbf{p}-\mathbf{p}_i^{\,n}\rVert^2$ — a rotation
    preserves length, and that is the cheapest check on the whole construction.};
\end{scope}

\end{tikzpicture}
\end{document}
